% Generated draft: requires mathematical and source review.
\documentclass[11pt]{article}
\usepackage[margin=25mm]{geometry}
\usepackage{fontspec}
\setmainfont{Latin Modern Roman}
\usepackage{amsmath,amssymb,mathtools,unicode-math}
\setmathfont{Latin Modern Math}
\usepackage{enumitem,needspace,xurl,hyperref}
\setlength{\parindent}{0pt}
\setlength{\parskip}{5pt}
\setlength{\emergencystretch}{3em}
\newcommand{\N}{\mathbb{N}}
\newcommand{\Z}{\mathbb{Z}}
\newcommand{\Q}{\mathbb{Q}}
\newcommand{\R}{\mathbb{R}}
\newcommand{\C}{\mathbb{C}}
\newcommand{\F}{\mathbb{F}}
\newcommand{\A}{\mathbb{A}}
\newcommand{\PP}{\mathbb{P}}
\newcommand{\E}{\mathbb{E}}
\newcommand{\eps}{\varepsilon}
\newcommand{\dd}{\,\mathrm{d}}
\begin{document}
\section*{UnsolvedMath ID: 20001928}
\textbf{Source problem code:} catalogue entry
\subsection{Definitions and mathematical statement}
Bangert's Question 10.1.1 asks whether a closed surface of genus at least two admits a Riemannian metric with completely integrable geodesic flow. The source distinguishes independence of the additional integral on an open everywhere dense set from independence on a set additionally having full measure.

In the smooth interpretation used by the aid, let \(H(x,p)=\tfrac12\|p\|_g^2\) on \(T^*M\), and let \(F\) be a smooth first integral, \(\{F,H\}=0\). On the unit energy shell \(S=H^{-1}(1/2)\), write \(f=F|_S\) and
\[
 \mathcal R=\{v\in S:df_v\ne0\},\qquad
 \mathcal C=\operatorname{Crit}(f)=\{v\in S:df_v=0\}.
\]
Independence of \(dF,dH\) is equivalent to this nonvanishing of the restricted differential. The point-critical set \(\mathcal C\) is not the larger inverse image of all critical values. Let \(\phi^t\) be the geodesic flow and let \(m_L\) be normalized Liouville measure on \(S\).

The paired aid reports that every ergodic \(\phi^1\)-invariant probability measure \(\mu\) with \(\mu(\mathcal C)=0\) has zero metric entropy. It consequently reports
\[
 h_{\mathrm{top}}(\phi^1|_S)=
 h_{\mathrm{top}}(\phi^1|_{\mathcal C}).
\]
For every metric on a closed genus-at-least-two surface, the aid uses positive topological entropy. Thus an integral whose regular set \(\mathcal R\) is open and dense must have a closed nowhere-dense point-critical set carrying all that topological entropy. In the full-measure version, \(m_L(\mathcal C)=0\) as well, and the reported conclusion is
\[
 h_{m_L}(\phi^1)=0.
\]
Positive topological entropy and zero Liouville entropy are not contradictory: other invariant measures may carry entropy on \(\mathcal C\). These are necessary conditions, not an existence or nonexistence theorem for complete smooth integrability.

The original context attributes to Ma\~n\'e [Ma\~n\'e1996] a generic Anosov-or-zero-Liouville-exponents statement for Hamiltonian flows on surfaces. The topology, energy-level conventions and exact genericity statement are omitted; it is retained as source context rather than a fully specified theorem here. The text ends by introducing an ``easier version'' of Bangert's question, but that subsequent question is not included and is not guessed.
\subsection{Short English statement}
The source asks for smooth complete integrability on a higher-genus surface and distinguishes dense-open from full-measure regularity. The aid reports that any candidate's entropy must be concentrated on the actual critical points of its integral. In the full-measure version that entropy-carrying set is Liouville-null, while Liouville entropy is zero; the existence question remains unresolved.
\subsection{Sources}
\begin{itemize}
\item Local source record: catalogue entry, catalogue entry; canonical \texttt{aim-geometry-notes.json}, \texttt{notes[265]}, source item 10.1, Question 10.1.1 attributed to Bangert. Title: ``A point-critical entropy-carrier obstruction to smooth integrability.''
\item Keith Burns and Vladimir S. Matveev, \emph{Open problems and questions about geodesics}, \emph{Ergodic Theory and Dynamical Systems} 41 (2021), 641--684, DOI \texttt{10.1017/etds.2019.73}; originating in the AIM ``Geodesics'' workshop. Supplied PDF reference \url{https://aimath.org/pastworkshops/geodesicsproblems.pdf}; this draft uses the supplied text without accessing a PDF.
\item Original reference [Ma\~n\'e1996], with the name typographically damaged in extraction. No full entry or exact genericity hypotheses are supplied. The point-critical localization and positive-entropy obstruction are claims of research attempt 1.
\item Registry dated 14 August 2026: ``partially\_solved,'' ``PARTIAL-PROGRESS,'' statement status ``exact''; attempt partial, research-run status open, rigor ``proved,'' novelty confidence low. The machine-generated, non-peer-reviewed aid expressly requires independent verification before citation.
\end{itemize}
\end{document}
